\chapter{Introduction}
\label{ch:introduction}

\begin{chaptersummary}
This chapter motivates the thesis from what tool-using language-model agents now do: they act inside
applications on a person's behalf, and they increasingly run on open models that an organisation can
host itself. It identifies five gaps that remain open when such agents work over many steps, states the
research problem, the aim, four objectives and six research questions, lists the contributions, and
outlines the thesis.
\end{chaptersummary}

\section{Background and Motivation}
\label{sec:intro-background}

Large language models (LLMs) are no longer used only to generate text. In an \emph{agentic} deployment
the model decomposes a request into steps, selects application programming interfaces (APIs), calls
them, reads the results, and decides what to do next \citep{yao2023react,schick2023toolformer}. Public
catalogues expose tens of thousands of callable endpoints; ToolBench alone indexes 16{,}464 RapidAPI
endpoints \citep{qin2024toolllm}. Multi-agent frameworks divide the work among several specialised
agents \citep{wu2024autogen,hong2024metagpt,qian2024chatdev,li2023camel}, and surveys identify their
coordination as an open systems problem \citep{wang2024survey,guo2024mas}.

The setting has changed in two ways that matter for this thesis. First, agents now \emph{act}: they send
money, delete files, post messages and change records inside applications, so a wrong step has a cost
that a wrong sentence does not. AppWorld \citep{trivedi2024appworld} makes this measurable with nine
everyday applications, 457 APIs and 750 tasks graded by unit tests on the final state of the world,
including tests that fail when an agent changes something it was not asked to change. Second,
reliability compounds: a chain of ten steps that are each 95\,\% reliable succeeds only about 60\,\% of
the time, with no single step looking faulty. On the AppWorld leaderboard (November 2026) the strongest
systems complete 91--94\,\% of the normal test tasks, but only with frontier models; the best entry that
uses a model of at most 14 billion parameters reaches 86.9\,\%, and no entry reports cost, collateral
damage or robustness to manipulated tool output.

Each step of such an agent runs the same four operations: (i) find the APIs relevant to this step, (ii)
decide which model should write the step, (iii) decide whether the proposed action can be trusted
before it changes the world, and (iv) decide how much of the growing history to keep. The literature
studies these operations separately and mostly per step:
\begin{itemize}[leftmargin=*]
\item Tool retrievers rank APIs from the current query only, ignoring which applications and calls the
episode has already committed to \citep{qin2024toolllm,zheng2024toolrerank}.
\item LLM routers choose a model from the current query and a cost budget, without regard for how risky
the episode has become \citep{ong2025routellm,panda2025pilot,yue2025masrouter}.
\item Consensus methods count votes over generated text, so several samples that agree on the same
mistake out-vote a correct one \citep{wang2023selfconsistency,du2024debate}; no method certifies an
\emph{action} before it executes.
\item Prompt compressors shrink context to a ratio fixed in advance, with no guarantee that the facts the
next step needs survive \citep{jiang2023llmlingua,li2023selectivecontext}.
\end{itemize}

The motivation of this research is that the signals each operation needs already exist inside the
episode: the applications the agent has used, the state changes candidate actions would cause, the
episode's running error rate, and the facts that earlier steps established. A system that carries
these signals across steps, and protects each with an explicit guarantee or an abstention rule, should
be more reliable and cheaper than four independently optimised stages, and should make that
reliability achievable with open models on modest hardware.

The same agent design addresses a public problem in India. Kisan Call Centres answered 33.53 lakh
farmer calls in 2022--23, and the Ministry of Agriculture and Farmers Welfare has announced chatbot and
AI-based decision support for them \citep{loksabha2023kcc}; Kisan e-Mitra already answers about 13{,}000
scheme queries a day \citep{wadhwani2024kisanemitra}. A deployable advisory agent must ground answers in
government data, leave an audit trail, and know which answers (a pesticide dose, an entitlement) it may
recommend but not decide.

\section{Problem Statement}
\label{sec:intro-problem}

Given an episode of steps over a large API catalogue and a choice of open models of different sizes,
current agents (a) retrieve APIs without episode memory, (b) choose models without regard to risk or
cost, (c) execute state-changing actions without verifying them and without any bound on how often they
are wrong, and (d) truncate history with no fidelity guarantee. The research problem is to design a
closed-loop agent architecture in which state produced at step $t$ improves retrieval, routing,
verification and context management at step $t+1$; to give the decision to act a certified error bound
and an explicit authority model; and to show, by controlled, leak-free experiments on simulated agents,
on a public-problem dataset and on real open-weight agents, which part of the loop contributes what.

\section{Research Aim and Objectives}
\label{sec:intro-objectives}

\textbf{Aim.} To design, implement and evaluate the Adaptive Context Reasoning System (ACRS), a
closed-loop architecture for tool-using LLM agents in which session-aware retrieval, risk-weighted
routing, certified consensus and fidelity-controlled pruning share one episode state, so that open
models on free or low-cost GPUs act reliably, cheaply and within stated error bounds.

\textbf{Objectives.}
\begin{enumerate}[label=\textbf{O\arabic*}, leftmargin=*]
\item \textbf{SATR -- Session-Adaptive Tool Scoring.} To design a training-free retriever that scores
APIs by their utility to the current step given what the episode has already used, so a small model
sees the right documentation within a fixed context budget (Chapter~\ref{ch:satr}).
\item \textbf{APRR -- Risk-Weighted Policy Routing.} To design a contextual-bandit router that assigns
each step to a small or a large model, whose exploration changes sign with the episode's risk, with an
escalation policy and a regret guarantee (Chapter~\ref{ch:aprr}).
\item \textbf{MNCD -- Byzantine-Tolerant Consensus and Certified Execution.} To design a consensus
mechanism that certifies an answer or a state-changing action only when independent candidates agree on
grounded evidence or on the state change they cause, abstains otherwise, tolerates up to $f$ faulty
participants, and sets its certification threshold with a finite-sample error guarantee, under an
explicit policy-as-code authority model (Chapters~\ref{ch:mncd} and~\ref{ch:governance}).
\item \textbf{FCNP -- PID-Controlled Provenance Pruning.} To design a context governor whose
proportional--integral controller removes as much history as a fidelity floor allows while never
evicting the facts and schemas a committed call needs (Chapter~\ref{ch:fcnp}).
\end{enumerate}

\section{Research Questions}
\label{sec:intro-rq}

\begin{description}[leftmargin=1.3cm, style=nextline]
\item[RQ1] Does episode state improve API retrieval over stateless retrieval once training-label leakage
is removed, and where does it stop helping?
\item[RQ2] Does risk-weighted routing lower cost and regret without raising errors?
\item[RQ3] Does consensus over grounded evidence or over executed outcomes prevent wrong and damaging
actions, including under prompt injection and collusion, and can its error rate be certified?
\item[RQ4] Does controlled pruning keep the facts an agent needs at a lower token cost than fixed-ratio
and recency baselines?
\item[RQ5] Does the closed loop outperform the four mechanisms deployed separately, on simulated agents,
on real farmer queries, and on real open-weight agents?
\item[RQ6] How close can open models on free or low-cost GPUs come to frontier-model agents when the loop
is in place, and at what cost and collateral damage?
\end{description}

\section{Contributions of the Thesis}
\label{sec:intro-contrib}

\begin{enumerate}[label=\textbf{C\arabic*}, leftmargin=*]
\item \textbf{A leakage audit and leak-free protocol} for session-level tool retrieval on ToolBench G1:
94 of the 100 test gold tool sets appear verbatim in the training trajectories, inflating NDCG@5 from
0.485 to 0.576 (Chapter~\ref{ch:methodology}).
\item \textbf{SATR}, a training-free session-aware retriever: +12.4\,\% held-out NDCG@5 over BM25 on
ToolBench (95\,\% CI of the difference +0.036 to +0.094), with a measured scope limit on topic-switching
queues (Chapters~\ref{ch:satr} and~\ref{ch:governance}).
\item \textbf{APRR}, a risk-signed contextual-bandit router with escalation and a regret bound on
optimistic rounds; lowest regret of six routers on ToolBench (Chapter~\ref{ch:aprr}).
\item \textbf{MNCD with certified execution}: an evidence-weighted Byzantine quorum (accepted errors
24.0\,\% to 0.6\,\% on ToolBench), extended to outcome-level consensus over executed candidate
programs, a Learn-then-Test certification threshold \citep{angelopoulos2025ltt}, and a policy-as-code
authority stage that is deny-by-default, tested both ways and audited in a hash chain
(Chapters~\ref{ch:mncd} and~\ref{ch:governance}).
\item \textbf{FCNP}, a typed-block pruner with a PI fidelity controller: fidelity floor met on 99.7\,\% of
turns while removing 47\,\% of tokens (Chapter~\ref{ch:fcnp}).
\item \textbf{KCC-Sessions}, a public-problem benchmark built from real Kisan Call Centre queries without
inventing labels, on which the integrated loop cuts executed errors from 11.3\,\% to 0.51\,\% and the
success of injection attacks from 27.8\,\% to 0.86\,\% (Chapter~\ref{ch:governance}).
\item \textbf{ACRS on AppWorld}: the four objectives specialised to an interactive environment with real
open-weight agents, evaluated with the official scorer, collateral-damage tests and paired significance
tests, together with an open, reproducible harness for free and low-cost GPUs
(Chapter~\ref{ch:appworld}).
\end{enumerate}

Figure~\ref{fig:contrib-map} maps each research gap to the objective that addresses it, the contribution
it produces and the chapter that holds the evidence.

\begin{figure}[htbp]
\centering
\includegraphics[width=\textwidth]{figures/01_contributions}
\caption{Research gaps, objectives, contributions and evidence}
\label{fig:contrib-map}
\end{figure}

\section{Organisation of the Thesis}
\label{sec:intro-organisation}

Chapter~\ref{ch:literature} reviews tool learning, LLM routing, multi-agent consensus, context
compression and agent safety, and derives the research gaps. Chapter~\ref{ch:methodology} describes the
design method, the ACRS loop, the three evaluation layers (ToolBench, KCC-Sessions and AppWorld), the
metrics and the statistical analysis. Chapters~\ref{ch:satr} to~\ref{ch:fcnp} present the four
objectives, each with its formulation, algorithm, formal properties, experiments and its specialisation
to AppWorld. Chapter~\ref{ch:integrated} evaluates the closed loop on ToolBench and a live AGMARKNET
case. Chapter~\ref{ch:governance} introduces authority-bounded execution and the Kisan Call Centre
evaluation. Chapter~\ref{ch:appworld} evaluates ACRS with real open-weight agents on AppWorld.
Chapter~\ref{ch:conclusion} answers the research questions, states the limitations and outlines future
work.
